\documentclass[10pt,a4paper]{amsart}
\usepackage{amsmath,amssymb,amsthm,graphicx}
\usepackage{verbatim}
\newenvironment{reference}{\comment}{\endcomment}
%\newenvironment{reference}{}{}
\newenvironment{slogan}{\comment}{\endcomment}
\newenvironment{history}{\comment}{\endcomment}

% For commutative diagrams we use Xy-pic
\usepackage[all]{xy}

% We use 2cell for 2-commutative diagrams.
\xyoption{2cell}
% Theorem environments.
%
\theoremstyle{plain}
\newtheorem{theorem}[subsection]{Theorem}
\newtheorem{proposition}[subsection]{Proposition}
\newtheorem{lemma}[subsection]{Lemma}

\theoremstyle{definition}
\newtheorem{definition}[subsection]{Definition}
\newtheorem{example}[subsection]{Example}
\newtheorem{exercise}[subsection]{Exercise}
\newtheorem{situation}[subsection]{Situation}

\theoremstyle{remark}
\newtheorem{remark}[subsection]{Remark}
\newtheorem{remarks}[subsection]{Remarks}

\numberwithin{equation}{subsection}

% Macros
%
\def\lim{\mathop{\mathrm{lim}}\nolimits}
\def\colim{\mathop{\mathrm{colim}}\nolimits}
\def\Spec{\mathop{\mathrm{Spec}}}
\def\Hom{\mathop{\mathrm{Hom}}\nolimits}
\def\Ext{\mathop{\mathrm{Ext}}\nolimits}
\def\SheafHom{\mathop{\mathcal{H}\!\mathit{om}}\nolimits}
\def\SheafExt{\mathop{\mathcal{E}\!\mathit{xt}}\nolimits}
\def\Sch{\mathit{Sch}}
\def\Mor{\mathop{\mathrm{Mor}}\nolimits}
\def\Ob{\mathop{\mathrm{Ob}}\nolimits}
\def\Sh{\mathop{\mathit{Sh}}\nolimits}
\def\NL{\mathop{N\!L}\nolimits}
\def\CH{\mathop{\mathrm{CH}}\nolimits}
\def\proetale{{pro\text{-}\acute{e}tale}}
\def\etale{{\acute{e}tale}}
\def\QCoh{\mathit{QCoh}}
\def\Ker{\mathop{\mathrm{Ker}}}
\def\Im{\mathop{\mathrm{Im}}}
\def\Coker{\mathop{\mathrm{Coker}}}
\def\Coim{\mathop{\mathrm{Coim}}}

% Boxtimes
%
\DeclareMathSymbol{\boxtimes}{\mathbin}{AMSa}{"02}

%
% Macros for moduli stacks/spaces
%
\def\QCohstack{\mathcal{QC}\!\mathit{oh}}
\def\Cohstack{\mathcal{C}\!\mathit{oh}}
\def\Spacesstack{\mathcal{S}\!\mathit{paces}}
\def\Quotfunctor{\mathrm{Quot}}
\def\Hilbfunctor{\mathrm{Hilb}}
\def\Curvesstack{\mathcal{C}\!\mathit{urves}}
\def\Polarizedstack{\mathcal{P}\!\mathit{olarized}}
\def\Complexesstack{\mathcal{C}\!\mathit{omplexes}}
% \Pic is the operator that assigns to X its picard group, usage \Pic(X)
% \Picardstack_{X/B} denotes the Picard stack of X over B
% \Picardfunctor_{X/B} denotes the Picard functor of X over B
\def\Pic{\mathop{\mathrm{Pic}}\nolimits}
\def\Picardstack{\mathcal{P}\!\mathit{ic}}
\def\Picardfunctor{\mathrm{Pic}}
\def\Deformationcategory{\mathcal{D}\!\mathit{ef}}
\usepackage{cleveref}
\crefname{originallemma}{original lemma}{original lemmas}
\Crefname{originallemma}{Original lemma}{Original lemmas}
\makeatletter
\newlabel{algebra-lemma-Artin-Rees}{{00IN}{source}{}{sourcepointer-1}{}}
\newlabel{algebra-lemma-Artin-Rees@cref}{{[originallemma][0][]00IN}{[1][0][]source}{}{}{}}
\makeatother
\crefname{originallemma}{original lemma}{original lemmas}
\Crefname{originallemma}{Original lemma}{Original lemmas}
\makeatletter
\newlabel{algebra-lemma-completion-tensor}{{00MA}{source}{}{sourcepointer-2}{}}
\newlabel{algebra-lemma-completion-tensor@cref}{{[originallemma][0][]00MA}{[1][0][]source}{}{}{}}
\makeatother
\crefname{originalsection}{original section}{original sections}
\Crefname{originalsection}{Original section}{Original sections}
\makeatletter
\newlabel{algebra-section-netherlander}{{00S0}{source}{}{sourcepointer-3}{}}
\newlabel{algebra-section-netherlander@cref}{{[originalsection][0][]00S0}{[1][0][]source}{}{}{}}
\makeatother
\crefname{originalsection}{original section}{original sections}
\Crefname{originalsection}{Original section}{Original sections}
\makeatletter
\newlabel{derived-section-derived-limit}{{08TB}{source}{}{sourcepointer-4}{}}
\newlabel{derived-section-derived-limit@cref}{{[originalsection][0][]08TB}{[1][0][]source}{}{}{}}
\makeatother
\crefname{originalequation}{original equation}{original equations}
\Crefname{originalequation}{Original equation}{Original equations}
\makeatletter
\newlabel{equation-C-prime}{{0AL4}{source}{}{sourcepointer-5}{}}
\newlabel{equation-C-prime@cref}{{[originalequation][0][]0AL4}{[1][0][]source}{}{}{}}
\makeatother
\crefname{originalreference}{original reference}{original references}
\Crefname{originalreference}{Original reference}{Original references}
\makeatletter
\newlabel{lemma-topologically-finite-type-Noetherian}{{0AJQ}{source}{}{sourcepointer-6}{}}
\newlabel{lemma-topologically-finite-type-Noetherian@cref}{{[originalreference][0][]0AJQ}{[1][0][]source}{}{}{}}
\makeatother
\crefname{originallemma}{original lemma}{original lemmas}
\Crefname{originallemma}{Original lemma}{Original lemmas}
\makeatletter
\newlabel{more-algebra-lemma-compute-Rlim-modules}{{091D}{source}{}{sourcepointer-7}{}}
\newlabel{more-algebra-lemma-compute-Rlim-modules@cref}{{[originallemma][0][]091D}{[1][0][]source}{}{}{}}
\makeatother
\crefname{originalremark}{original remark}{original remarks}
\Crefname{originalremark}{Original remark}{Original remarks}
\makeatletter
\newlabel{more-algebra-remark-how-unique}{{07KZ}{source}{}{sourcepointer-8}{}}
\newlabel{more-algebra-remark-how-unique@cref}{{[originalremark][0][]07KZ}{[1][0][]source}{}{}{}}
\makeatother
\crefname{originalsection}{original section}{original sections}
\Crefname{originalsection}{Original section}{Original sections}
\makeatletter
\newlabel{section-two-categories}{{0AL2}{source}{}{sourcepointer-9}{}}
\newlabel{section-two-categories@cref}{{[originalsection][0][]0AL2}{[1][0][]source}{}{}{}}
\makeatother
\title{The completed naively trivialized cotangent complex as the derived limit of its reductions}
\author{The Stacks Project Authors}
\date{stacks-project:restricted.tex:local-snapshot-20261012; GNU FDL 1.2 or later}
\begin{document}
\maketitle
\noindent\textit{Editorial note.} All delivered mathematical text is the author's own, copied from the pinned Stacks Project chapter file restricted.tex; no statement or proof step was rewritten, reordered or removed.\par\medskip
\noindent\textit{Editorial note.} The counted proof is the author's complete argument: the source gives the full Artin-Rees comparison chain and the completeness argument for (2); the compatibility of the canonical quotient maps with differentials is stated as a direct check ('a moment's thought')\par\medskip
\noindent\textit{Editorial note.} The label lemma-NL-is-limit names the counted statement, and the proof environment immediately following it in the source is delivered verbatim as its proof.\par\medskip
\noindent\textit{Editorial note.} The counted statement is carried by the author's own lemma environment in the chapter source.\par\medskip
\noindent\textit{Allowed standard prerequisites (editorial).}\begin{itemize}
\item the naively trivialized cotangent complex NL of a presentation and its behaviour under base change
\item the two categories of the chapter and the rings B\_n = B/I\textasciicircum{}nB, A\_n
\item Artin-Rees for Noetherian rings
\item homotopy limits Rlim of inverse systems of complexes and pro-isomorphic systems
\end{itemize}\medskip
\ifdefined\hypertarget\hypertarget{sourcepointer-1}{}\fi
\noindent\textit{Source reference.} Original lemma 00IN: sources/stacks-chow/tags.txt, source lines 685--685. Prerequisite pointer invoked at proof line(s) 636; it names a standard fact of a sibling chapter that the counted argument uses, and every step of the counted proof built on it is retained verbatim in the delivered body.\par\medskip
\ifdefined\hypertarget\hypertarget{sourcepointer-2}{}\fi
\noindent\textit{Source reference.} Original lemma 00MA: sources/stacks-chow/tags.txt, source lines 816--816. Prerequisite pointer invoked at proof line(s) 662; it names a standard fact of a sibling chapter that the counted argument uses, and every step of the counted proof built on it is retained verbatim in the delivered body.\par\medskip
\ifdefined\hypertarget\hypertarget{sourcepointer-3}{}\fi
\noindent\textit{Source reference.} Original section 00S0: sources/stacks-chow/tags.txt, source lines 1022--1022. Prerequisite pointer invoked at proof line(s) 612; it names a standard fact of a sibling chapter that the counted argument uses, and every step of the counted proof built on it is retained verbatim in the delivered body.\par\medskip
\ifdefined\hypertarget\hypertarget{sourcepointer-4}{}\fi
\noindent\textit{Source reference.} Original section 08TB: sources/stacks-chow/tags.txt, source lines 11159--11159. Prerequisite pointer invoked at proof line(s) 620; it names a standard fact of a sibling chapter that the counted argument uses, and every step of the counted proof built on it is retained verbatim in the delivered body.\par\medskip
\ifdefined\hypertarget\hypertarget{sourcepointer-5}{}\fi
\noindent\textit{Source reference.} Original equation 0AL4: sources/stacks-restricted/pinned-original-source.tex, source lines 326--326. Prerequisite pointer invoked at proof line(s) 598; it names a standard fact of the same chapter that the counted argument uses, and every step of the counted proof built on it is retained verbatim in the delivered body.\par\medskip
\ifdefined\hypertarget\hypertarget{sourcepointer-6}{}\fi
\noindent\textit{Source reference.} Original reference 0AJQ: sources/stacks-restricted/pinned-original-source.tex, source lines 395--395. Prerequisite pointer invoked at proof line(s) 638; it names a standard fact of the same chapter that the counted argument uses, and every step of the counted proof built on it is retained verbatim in the delivered body.\par\medskip
\ifdefined\hypertarget\hypertarget{sourcepointer-7}{}\fi
\noindent\textit{Source reference.} Original lemma 091D: sources/stacks-chow/tags.txt, source lines 11441--11441. Prerequisite pointer invoked at proof line(s) 664; it names a standard fact of a sibling chapter that the counted argument uses, and every step of the counted proof built on it is retained verbatim in the delivered body.\par\medskip
\ifdefined\hypertarget\hypertarget{sourcepointer-8}{}\fi
\noindent\textit{Source reference.} Original remark 07KZ: sources/stacks-chow/tags.txt, source lines 9674--9674. Prerequisite pointer invoked at proof line(s) 665; it names a standard fact of a sibling chapter that the counted argument uses, and every step of the counted proof built on it is retained verbatim in the delivered body.\par\medskip
\ifdefined\hypertarget\hypertarget{sourcepointer-9}{}\fi
\noindent\textit{Source reference.} Original section 0AL2: sources/stacks-restricted/pinned-original-source.tex, source lines 303--303. Prerequisite pointer invoked at proof line(s) 605; it names a standard fact of the same chapter that the counted argument uses, and every step of the counted proof built on it is retained verbatim in the delivered body.\par\medskip
\begin{lemma}
\label{lemma-NL-is-limit}
Let $A$ be a Noetherian ring and let $I \subset A$ be an ideal.
Let $B$ be an object of (\ref{equation-C-prime}). Then
\begin{enumerate}
\item the pro-objects
$\{\NL_{B/A}^\wedge \otimes_B B/I^nB\}$ and $\{\NL_{B_n/A_n}\}$
of $D(B)$ are strictly isomorphic (see proof for elucidation),
\item $\NL_{B/A}^\wedge = R\lim \NL_{B_n/A_n}$ in $D(B)$.
\end{enumerate}
Here $B_n$ and $A_n$ are as in Section \ref{section-two-categories}.
\end{lemma}

\begin{proof}
The statement means the following: for every $n$ we have a well
defined complex $\NL_{B_n/A_n}$ of $B_n$-modules and we have
transition maps $\NL_{B_{n + 1}/A_{n + 1}} \to \NL_{B_n/A_n}$.
See Algebra, Section \ref{algebra-section-netherlander}.
Thus we can consider
$$
\ldots \to \NL_{B_3/A_3} \to \NL_{B_2/A_2} \to \NL_{B_1/A_1}
$$
as an inverse system of complexes of $B$-modules and a fortiori as
an inverse system in $D(B)$. Furthermore $R\lim \NL_{B_n/A_n}$
is a homotopy limit of this inverse system, see
Derived Categories, Section \ref{derived-section-derived-limit}.

\medskip\noindent
Choose a presentation $B = A[x_1, \ldots, x_r]^\wedge / J$. This
defines presentations
$$
B_n = B/I^nB = A_n[x_1, \ldots, x_r]/J_n
$$
where
$$
J_n = JA_n[x_1, \ldots, x_r] =
J/(J \cap I^nA[x_1, \ldots, x_r]^\wedge)
$$
The two term complex $J_n/J_n^2 \longrightarrow \bigoplus B_n \text{d}x_i$
represents $\NL_{B_n/A_n}$, see
Algebra, Section \ref{algebra-section-netherlander}.
By Artin-Rees (Algebra, Lemma \ref{algebra-lemma-Artin-Rees})
in the Noetherian ring $A[x_1, \ldots, x_r]^\wedge$
(Lemma \ref{lemma-topologically-finite-type-Noetherian})
we find a $c \geq 0$ such that we have canonical surjections
$$
J/I^nJ \to J_n \to J/I^{n - c}J \to J_{n - c},\quad n \geq c
$$
for all $n \geq c$. A moment's thought shows that these maps are
compatible with differentials and we obtain maps of complexes
$$
\NL_{B/A}^\wedge \otimes_B B/I^nB \to
\NL_{B_n/A_n} \to
\NL_{B/A}^\wedge \otimes_B B/I^{n - c}B \to
\NL_{B_{n - c}/A_{n - c}}
$$
compatible with the transition maps of the inverse systems
$\{\NL_{B/A}^\wedge \otimes_B B/I^nB\}$ and $\{\NL_{B_n/A_n}\}$.
This proves part (1) of the lemma.

\medskip\noindent
By part (1) and since pro-isomorphic
systems have the same $R\lim$ in order to prove (2)
it suffices to show that $\NL_{B/A}^\wedge$ is
equal to $R\lim \NL_{B/A}^\wedge \otimes_B B/I^nB$.
However, $\NL_{B/A}^\wedge$ is a two term complex $M^\bullet$
of finite $B$-modules which are $I$-adically complete for example by
Algebra, Lemma \ref{algebra-lemma-completion-tensor}. Hence
$M^\bullet = \lim M^\bullet/I^nM^\bullet = R\lim M^\bullet/I^n M^\bullet$, see
More on Algebra, Lemma \ref{more-algebra-lemma-compute-Rlim-modules} and
Remark \ref{more-algebra-remark-how-unique}.
\end{proof}

\end{document}
